\section{Область применимости условия кривизны}\label{sec:structure}
\subsection{Обычная гладкость и модель через норму градиента}
Поскольку $f$ дважды непрерывно дифференцируема в окрестности
компактного $Q$, общая граница $L=\max_{x\in Q}\|\nabla^2f(x)\|$
конечна. Польза нового условия состоит в параметрах и локальной
модели, а не в охвате функций без конечной константы гладкости на
данном фиксированном $Q$. Кроме того,
\[
 G(x)\le D\|\nabla f(x)\|,
 \qquad \|\nabla^2f(x)\|\le L_0+(L_1D)\|\nabla f(x)\|.
\]
Значит, гладкость через зазор влечёт оценку через норму градиента на $Q$.
Обратный переход с нулевым постоянным слагаемым неверен:
для $f(t)=e^t$ на $[0,R]$ выполнено $|f''|=|f'|$, но $G(0)=0$
и $f''(0)=1$. Поэтому любая допустимая пара параметров через зазор
в этом примере имеет $L_0\ge1$. В общем случае
$L_0\ge\|\nabla^2f(x^*)\|$.

\begin{proposition}[Ограничение режима с нулевым постоянным слагаемым]
\label{prop:rigidity}
Пусть \eqref{eq:gap_smoothness} выполнено при $L_0=0$.
Если $\nabla f(z)=0$ хотя бы в одной точке $z\in Q$, то
$\nabla f=0$ на $Q$ и $f$ постоянна на $Q$.
Выпуклость или квазар-выпуклость $f$ не требуется.
\end{proposition}
\begin{proof}
Зафиксируем $y\in Q$ и положим $u(t)=\|\nabla f(z+t(y-z))\|$.
Основная теорема анализа и $G\le D\|\nabla f\|$ дают
\[
 u(t)\le L_1D\|y-z\|\int_0^t u(s)\,ds.
\]
По неравенству Гронуолла $u=0$. Точка $y$ произвольна.
\end{proof}
Следовательно, непостоянная функция в этом режиме не может иметь
минимизатор с нулевым градиентом. В частности, исключён внутренний
минимум на полноразмерном множестве. Это ограничивает область
применимости следствия~\ref{cor:pure_gap_linear}.

\subsection{Связь с известными условиями роста}
Для выпуклой $f$ обозначим
$D_f(y,x)=f(y)-f(x)-\langle\nabla f(x),y-x\rangle$.
При $L_0=0<L_1$ лемма~\ref{lem:transport} даёт для $0\le\alpha\le1$
\[
 D_f(x+\alpha d,x)\le\frac{M\alpha^2}{2}G(x),
 \qquad M=L_1D^2e^{L_1D^2}.
\]
Достаточно ограничить экспоненту в \eqref{eq:upper_model} числом
$e^{L_1D^2}$ и проинтегрировать $\alpha-t$.
Это strong $(M,1)$-growth из~\cite{wirth}.
Такое сведение консервативно: возникает экспоненциальная зависимость
от $L_1D^2$, тогда как следствие~\ref{cor:pure_gap_linear} использует
нижнюю границу шага порядка $\min\{1,(L_1D^2)^{-1}\}$.
Сравниваются именно выписанные способы преобразования констант;
это не нижняя граница возможностей всех других методов или доказательств.

\subsection{Многомерный класс с явными параметрами}
Пусть $a_i\in\mathbb R^n$, $y_i\ge0$, $B=\max_i\|a_i\|>0$.
Рассмотрим отрицательное логарифмическое правдоподобие пуассоновской
модели без слагаемых, зависящих только от данных:
\begin{equation}\label{eq:poisson}
 f(x)=\frac1m\sum_{i=1}^m\bigl(e^{a_i^Tx}-y_i a_i^Tx\bigr),
 \qquad 0\in Q.
\end{equation}
При целых $y_i$ это задача регрессии счётных данных; оценка ниже
верна и для неотрицательных вещественных откликов.

\begin{proposition}[Параметры из наблюдаемых данных]\label{prop:poisson}
На любом компактном выпуклом $Q$, содержащем ноль, функция
\eqref{eq:poisson} удовлетворяет \eqref{eq:gap_smoothness} при
\begin{equation}\label{eq:poisson_constants}
 L_0=B^2\left(1+\frac{1}{2m}\sum_i y_i^2\right),
 \qquad L_1=B^2,\qquad \gamma=1.
\end{equation}
Минимизатор и значение минимума для этих констант не нужны.
Если $Q=\{x:\|x\|\le R\}$, то допустимы параметры модели через
норму градиента $K_0=L_0$, $K_1=RL_1$ и общая граница гессиана
\begin{equation}\label{eq:poisson_uniform}
 L=e^{BR}\|A\|_2^2/m,
\end{equation}
где строки $A$ равны $a_i^T$. Утверждается допустимость этих
границ, а не их минимальность.
\end{proposition}
\begin{proof}
При $y\ge0$ и любом $t\in\mathbb R$ верно
\begin{equation}\label{eq:poisson_scalar}
 e^t\le1+y^2/2+t(e^t-y).
\end{equation}
Если $t\le0$, то $(1-t)e^t+yt\le1$.
Если $t\ge0$, производная $(1-t)e^t$ равна $-te^t\le-t$,
поэтому $(1-t)e^t+yt\le1-t^2/2+yt\le1+y^2/2$.
При этом $\|\nabla^2f(x)\|\le B^2m^{-1}\sum_i e^{a_i^Tx}$.
Суммирование \eqref{eq:poisson_scalar} даёт
\[
 \|\nabla^2f(x)\|\le L_0+L_1\langle\nabla f(x),x\rangle
 \le L_0+L_1G(x),
\]
поскольку $0\in Q$.
На шаре $\langle\nabla f(x),x\rangle\le R\|\nabla f(x)\|$.
Наконец, из $a_i^Tx\le BR$ следует
$\nabla^2f(x)\preceq e^{BR}A^TA/m$.
\end{proof}
При ненулевом градиенте линейный оракул шара имеет вид
$s=-R\nabla f(x)/\|\nabla f(x)\|$ и в общем случае отличается от
минимизатора \eqref{eq:poisson}. Эти задачи не обладают вырожденным
свойством скалярных примеров ниже, где оракул сразу находит минимум.
Условия и параметры доказаны на всём $Q$; численные опыты лишь
проверяют их следствия в выбранных точках.
